% 第4章；原PDF物理页108--128，原印刷页101--121。
\SourcePage{108}
\par\medskip\noindent\textbf{4.0.} 本章及下一章，除讨论\term{vanishing-cycles}{消失轮换}的4.4节外，我们只考虑域 $k$ 上分离\term{finite-type}{有限型}的\term{scheme}{概形}，并且只考虑\term{self-dual}{自对偶}\term{perversity}{反常度} $p_{1/2}$。我们将在 $\D_c^b(X,\overline\Q_\ell)$ 中工作（2.2.12）。为了使这个范畴有定义并满足通常的运算形式，我们必须对域 $k$ 加以限制：对 $k$ 的每个有限\term{extension}{扩张} $k'$，各群 $H^i(\Gal(\bar k/k'),\Z/\ell)$ 均应有限（参见2.2.12）；例如，有限域和代数闭域都满足这个条件。

这里的大多数结果对 $\Q_\ell$、$\Z_\ell$ 或 $\Z/\ell^n$ 系数也成立，尤其适用于 $\Z/\ell$-\term{cohomology}{上同调}。也可以将 $\Q_\ell$ 换成它的有限\term{extension}{扩张} $E_\lambda$，将 $\Z_\ell$ 换成相应的赋值环 $V_\lambda$，将 $\Z/\ell^n$ 换成 $V_\lambda$ 的有限商 $R$。若把所有这些变式一一写出，文字就难以阅读了。因此，我们选择只讨论其中一种，即 $\overline\Q_\ell$ 系数的情形。证明常常归结为 $\Z/\ell\Z$-模\term{sheaf}{层}的情形，而这种归结往往留给读者。作此归结时，应记住：限制标量函子 $T$（从 $V_\lambda$ 到 $\Z_\ell$），以及从 $V_\lambda$ 到其有限商的\term{extension}{扩张}标量函子，例如从 $\Z_\ell$ 到 $\Z/\ell^n$ 的函子，都是\term{conservative}{保守函子}，并且 $K\in{}^p\D^{\le0}$ 当且仅当 $TK\in{}^p\D^{\le0}$。另见2.2.17、2.2.18。

读者若要把这里的结果用于 $\overline\Q_\ell$ 以外的系数，须注意以下困难。
\begin{enumerate}[label=(\alph*)]
\item 在 $\Z_\ell$-\term{cohomology}{上同调}中，\term{self-dual}{自对偶}\term{perversity}{反常度} $p_{1/2}$ 产生两个 \term{t-structure}{$t$-结构} $[p_{1/2}]$ 和 $[p_{1/2}^{+}]$（3.3），它们由\term{duality}{对偶}交换。
\item 定理4.3.1在 $\Z_\ell$-\term{cohomology}{上同调}中不成立，即使 $X$ 是一点也是如此：$[p_{1/2}]$-\term{perverse-sheaf}{反常层}范畴只有\term{noetherian}{诺特}性，$[p_{1/2}^{+}]$-\term{perverse-sheaf}{反常层}范畴只有\term{artinian}{阿廷}性。
\item 4.5节的讨论只有在系数为域时才有意义。
\item 当系数为域时，张量积具有更好的正合性；例如，4.2.8要求系数为域。
\end{enumerate}

\SourcePage{109}
回顾（参见2.2.11），对于\term{self-dual}{自对偶}\term{perversity}{反常度} $p_{1/2}$，$X$ 上的一个\term{perverse-sheaf}{反常层}是 $\D_c^b(X,\overline\Q_\ell)$ 中的对象 $K$，满足：对 $X$ 的每个点 $x$（不必为闭点），令
\[\dim(x)=\dim\overline{\{x\}}=\operatorname{trdeg}(k(x)/k),\]
并以 $i_x$ 表示 $x$ 到 $X$ 的包含映射，则
\begin{equation*}H^i i_x^*K=0\quad(i>-\dim(x)),\tag{4.0.1}\end{equation*}
\begin{equation*}H^i i_x^!K=0\quad(i<-\dim(x)).\tag{4.0.2}\end{equation*}
条件（4.0.1）刻画 ${}^p\D_c^{\le0}$，条件（4.0.2）刻画 ${}^p\D_c^{\ge0}$。条件（4.0.1）\term{equivalence}{等价}于：对每个 $i$ 都有
\begin{equation*}\dim\Supp H^iK\le-i.\tag{4.0.1$'$}\end{equation*}
特别地，$i>0$ 时 $H^iK=0$。

\noindent\textbf{例。} 若 $X$ \term{smooth}{光滑}且\term{pure}{纯} $d$ 维，则 $\overline\Q_\ell[d]$（更一般地，放在 $-d$ 次的\term{smooth}{光滑}\term{sheaf}{层}）是\term{perverse-sheaf}{反常层}；放在\term{cohomology}{上同调}次数 $0$ 的点\term{sheaf}{层}也是\term{perverse-sheaf}{反常层}（参见2.2.2）。

由于\term{perversity}{反常度} $p_{1/2}$ \term{self-dual}{自对偶}，Verdier\term{duality}{对偶}交换 ${}^p\D^{\ge0}$ 与 ${}^p\D^{\le0}$，并在\term{perverse-sheaf}{反常层}范畴上诱导一个反自\term{equivalence}{等价}。

\par\medskip\noindent\textbf{4.0.3. 记号。} 对 $K\in\D_c^b(X,\overline\Q_\ell)$ 及任一 $X$ 上\term{scheme}{概形} $u:Z\to X$，我们以 $H^i(Z,K)$ 表示\term{hypercohomology}{超上同调}群 $\mathbb H^i(Z,u^*K)$。

\section*{4.1. 仿射态射}
\addcontentsline{toc}{section}{4.1. 仿射态射}
以下定理是M. Artin的定理SGA 4 XIV 3.1的另一种表述。

\par\medskip\noindent\textbf{定理4.1.1。} 若 $f:X\to Y$ 为\term{affine}{仿射}态射，则函子
\[Rf_*:\D_c^b(X,\overline\Q_\ell)\longrightarrow\D_c^b(Y,\overline\Q_\ell)\]
是右 \term{t-exact}{$t$-正合}的。

可归结为证明 $\Z/\ell$-\term{cohomology}{上同调}中的相应命题（参见4.0）。

若\term{sheaf}{层} $F$ \term{constructible}{可构造}，则M. Artin赋予 $F$ 的整数 $d(F)$，就是使得 $F$ 作为\term{derived-category}{导出范畴}的对象属于 ${}^p\D_c^{\le d}$ 的最小整数 $d$。对\term{derived-category}{导出范畴}中的\term{constructible}{可构造}对象 $K$，
\[d(K):=\sup_i\bigl(i+d(H^iK)\bigr)\]
同样是使 $K\in{}^p\D_c^{\le d}$ 的最小整数 $d$。$K$ 属于 ${}^p\D_c^{\le0}$ 的充要条件，是各个 $H^iK[-i]$ 都属于这个范畴。
\SourcePage{110}
因此，要验证 $Rf_*$ 将 ${}^p\D_c^{\le0}(X)$ 映入 ${}^p\D_c^{\le0}(Y)$，只需验证：对任意满足 $d(F)\le d$ 的\term{sheaf}{层} $F$，都有
\[d(R^if_*F)\le d-i.\]
这正是SGA 4 XIV 3.1的结论。

应用Verdier\term{duality}{对偶}，定理4.1.1\term{duality}{对偶}为如下推论。

\par\medskip\noindent\textbf{推论4.1.2。} 在相同假设下，$Rf_!$ 是左 \term{t-exact}{$t$-正合}的。

我们以 ${}^pf_*$ 和 ${}^pf_!$ 分别表示从 $X$ 上的\term{perverse-sheaf}{反常层}到 $Y$ 上\term{perverse-sheaf}{反常层}的函子 ${}^pH^0Rf_*$ 和 ${}^pH^0Rf_!$。前者右正合，后者左正合。对两个\term{affine}{仿射}态射的复合 $fg$，有
\[{}^p(fg)_*={}^pf_*\,{}^pg_*,\qquad{}^p(fg)_!={}^pf_!\,{}^pg_!\qquad\text{（1.3.17(iv)）}.\]
下文将 $Rf_*$、$Rf_!$ 简记为 $f_*$、$f_!$。

\par\medskip\noindent\textbf{推论4.1.3。} 若 $f:X\to Y$ 是拟有限\term{affine}{仿射}态射，则 $f_!$ 与 $f_*$ 均为 \term{t-exact}{$t$-正合}函子。

事实上，$f_*\,(:=Rf_*)$ 左 \term{t-exact}{$t$-正合}（2.2.5），也右 \term{t-exact}{$t$-正合}（4.1.7）。\footnote{译注：书末勘误第51条将此处的4.1.7更正为4.1.1。}$f_!$ 的情形与之\term{duality}{对偶}（2.2.5及4.1.2）。

设 $k$ 代数闭。取 $Y=\Spec(k)$，4.1.1就给出SGA 4 XIV 3.2中的推论。

\par\medskip\noindent\textbf{推论4.1.4。} 设 $X$ 为代数闭域 $k$ 上的\term{affine}{仿射}\term{scheme}{概形}，$F$ 为 $X$ 上任意\term{etale}{平展}\term{sheaf}{层}，则 $i>\dim X$ 时 $H^i(X,F)=0$。

更一般地，有如下推论。

\par\medskip\noindent\textbf{推论4.1.5。} 设 $k$ 代数闭，$K\in{}^p\D_c^{\le0}(X)$，$u:Z\to X$ 为源\term{scheme}{概形}\term{affine}{仿射}的拟有限态射，则 $i>0$ 时 $H^i(Z,K)=0$。

事实上，$u^*K\in{}^p\D_c^{\le0}(Z)$（2.2.5的类似结论），再对 $Z\to\Spec(k)$ 应用4.1.1即可。

这个推论有如下逆命题。

\par\medskip\noindent\textbf{逆命题4.1.6。} 设 $k$ 代数闭，$K\in\D_c^b(X)$。
\SourcePage{111}
若对 $X$ 的每个Zariski\term{affine}{仿射}开集 $V$ 都有 $H^i(V,K)=0$（$i>0$），则 $K\in{}^p\D_c^{\le0}(X)$。

我们先证明以下引理。

\par\medskip\noindent\textbf{引理4.1.7。} 设 $X$ 为代数闭域 $k$ 上的\term{smooth}{光滑}不可约\term{affine}{仿射}\term{scheme}{概形}，维数为 $d\ge1$，$F$ 为 $X$ 上的\term{smooth}{光滑}\term{sheaf}{层}。若 $F\ne0$，则存在形如 $U=X[1/f]$（$f\ne0$）的\term{affine}{仿射}开集，使 $H^d(U,F)\ne0$。

取 $f_1,\ldots,f_d\in\Gamma(X,\mathcal O_X)$，使方程 $f_i=0$ 定义的闭子\term{scheme}{概形} $S$ 有限且非空。若 $F\ne0$，由\term{pure}{纯}性得
\[H_S^{2d}(X,F)\simeq(F|S)(-d)\ne0.\]
\term{exact-sequence}{正合列}
\[H^{2d-1}(X-S,F)\longrightarrow H_S^{2d}(X,F)\longrightarrow H^{2d}(X,F)\]
表明 $H^{2d-1}(X-S,F)\ne0$。事实上，$2d>d$，故由4.1.4可知 $H^{2d}(X,F)=0$。以开集 $U_i:=X[1/f_i]$ 覆盖 $X-S$；对 $I\subset[1,d]$，令 $U_I=\bigcap_{i\in I}U_i$。各个 $U_I$ 都\term{affine}{仿射}，因此在Čech\term{spectral-sequence}{谱序列}
\[E_1^{pq}=\bigoplus_{|I|=p+1}H^q(U_I,F)\Longrightarrow H^{p+q}(X-S,F)\]
中，$q>d$ 时 $E_1^{pq}=0$。当 $p>d-1$ 时，$E_1^{pq}$ 也为零。只有 $E_1^{d-1,d}$ 对 $H^{2d-1}$ 有贡献，所以这一项非零，即
\[H^d\!\left(X\left[\left(\prod_i f_i\right)^{-1}\right],F\right)\ne0.\]

\par\medskip\noindent\textbf{4.1.8. 4.1.6的证明。} 取 $X$ 的一列递增闭集 $F_i$（$0\le i\le\dim X$），末项为 $X$，并令 $F_{-1}=\varnothing$。选择这列闭集，使 $W_i=F_i-F_{i-1}$ 赋予既约结构后\term{smooth}{光滑}且\term{pure}{纯} $i$ 维，且各个 $H^jK$ 在各 $W_i$ 上\term{smooth}{光滑}。我们对 $i$ 作下降归纳，证明在 $U_i=X-F_{i-1}$ 上，$K\in\D^{\le-i}$（这里使用自然 \term{t-structure}{$t$-结构}）。

当 $i=\dim X$ 时，归纳假设“$j>i$ 时 $K|U_j\in\D^{\le-j}$”为空。这个假设保证：在 $U_i$ 上，$\tau_{<-i}K\in{}^p\D^{\le0}$；而且，对 $j\ge-i$，$H^jK|U_i$ 是 $W_i$ 上的\term{smooth}{光滑}\term{sheaf}{层}的\term{extension-by-zero}{零延拓}。对 $U_i$ 的每个\term{affine}{仿射}开集 $V$，三角 $(\tau_{<-i}K,K,\tau_{\ge-i}K)$ 给出\term{cohomology}{上同调}长\term{exact-sequence}{正合列}
\SourcePage{112}
\[H^j(V,\tau_{<-i}K)\longrightarrow H^j(V,K)\longrightarrow H^j(V,\tau_{\ge-i}K)\longrightarrow\cdots.\]
由4.1.5，$j>0$ 时 $H^j(V,\tau_{<-i}K)=0$。再由假设 $H^j(V,K)=0$（$j>0$），得到当 $j>0$ 时
\[H^j(V,\tau_{\ge-i}K)=H^j(V\cap W_i,\tau_{\ge-i}K)=0.\]
由4.1.4，\term{spectral-sequence}{谱序列}
\[E_2^{pq}=H^p\bigl(V\cap W_i,H^q(\tau_{\ge-i}K)\bigr)\Longrightarrow H^{p+q}(V\cap W_i,\tau_{\ge-i}K)\]
的 $E_2$ 页在 $p>i$ 时为零。若 $j>q$ 时 $H^j(\tau_{\ge-i}K)=0$，便有
\[H^i\bigl(V\cap W_i,H^q(\tau_{\ge-i}K)\bigr)=H^{i+q}(V\cap W_i,\tau_{\ge-i}K),\]
当 $q>-i$ 时这个群为零。令 $L=H^q(\tau_{\ge-i}K)$，这是 $W_i$ 上的\term{smooth}{光滑}\term{sheaf}{层}。问题归结为由“对每个\term{affine}{仿射} $V\subset U_i$，$H^i(V\cap W_i,L)=0$”推出 $L=0$。$i=0$ 时显然。因而设 $i\ge1$。若 $L\ne0$，可取\term{affine}{仿射} $V_1\subset U_i$，使 $L|V_1\cap W_i\ne0$。对 $V_1\cap W_i$ 的不可约分支应用4.1.7，得到 $f_1\in\Gamma(V_1\cap W_i,\mathcal O)$，使
\[H^i(V_1\cap W_i[1/f_1],L)\ne0.\]
将 $f_1$ 提升为 $f\in\Gamma(V_1,\mathcal O)$，令 $V=V_1[1/f]$，就有 $H^i(V\cap W_i,L)\ne0$，矛盾。

\par\medskip\noindent\textbf{注4.1.9。} 当 $k=\C$ 时，4.1.1在复系数\term{cohomology}{上同调}中的类似结论也成立，并可利用比较定理由4.1.1推出。下面用\term{holonomic-module}{完整模}作一个局部的解释。

设 $X$ \term{smooth}{光滑}且\term{pure}{纯} $d$ 维。在 $X$ 上，带可积联络的\term{quasi-coherent}{拟凝聚}（代数）\term{sheaf}{层}，与作为 $\mathcal O$-模\term{quasi-coherent}{拟凝聚}的 $\mathscr D$-模，是同一概念（也可以说作为 $\mathscr D$-模\term{quasi-coherent}{拟凝聚}，两者\term{equivalence}{等价}）。每个这样的\term{sheaf}{层} $(V,\nabla)$ 都有相应的de Rham\term{complex}{复形} $\Omega^*(V)$；在 $X(\C)$ 上，还有解析de Rham\term{complex}{复形} $\Omega^*(V)^{\mathrm{an}}$，其各项为 $\Omega^i(V)\otimes_{\mathcal O}\mathcal O^{\mathrm{an}}$。已知若 $\mathscr D$-模 $V$ 完整，则 $\Omega^*(V)[d]^{\mathrm{an}}$ 为\term{perverse-sheaf}{反常层}，而且函子 $V\mapsto\Omega^*(V)[d]^{\mathrm{an}}$ 给出具有\term{regular-singularity}{正则奇点}（包括无穷远处）的完整 $\mathscr D$-模范畴与\term{perverse-sheaf}{反常层}范畴之间的\term{equivalence}{等价}。

若 $f:X\hookrightarrow Y$ 是到\term{smooth}{光滑}\term{pure}{纯} $d$ 维\term{scheme}{概形} $Y$ 的开浸入，
\SourcePage{113}
则在这一\term{equivalence}{等价}下，${}^pf_*$ 对应于\term{quasi-coherent}{拟凝聚}\term{sheaf}{层}的\term{direct-image}{正像}函子（联络随之得到）。若 $f$ 还\term{affine}{仿射}，则这一函子正合，正如4.1.3所预言的。

\par\medskip\noindent\textbf{推论4.1.10。} 设 $j:U\hookrightarrow X$ 为\term{affine}{仿射}开浸入，$i:D\hookrightarrow X$ 为其补闭集的包含映射（例如，$D$ 为Cartier除子而 $U=X-D$）。回顾：我们把 $Rj_*$、$Ri^!$ 简记为 $j_*$、$i^!$。
\begin{enumerate}[label=(\roman*)]
\item 若 $G$ 为 $U$ 上的\term{perverse-sheaf}{反常层}，则 $j_!G$ 与 $j_*G$ 均为\term{perverse-sheaf}{反常层}。
\item 若 $F$ 为 $X$ 上的\term{perverse-sheaf}{反常层}，则 $i^*F\in{}^p\D_c^{[-1,0]}$，$i^!F\in{}^p\D_c^{[0,1]}$。有\term{perverse-sheaf}{反常层}的\term{exact-sequence}{正合列}
\end{enumerate}
\begin{equation*}0\longrightarrow i_*{}^pH^{-1}i^*F\longrightarrow j_!j^*F\longrightarrow F\longrightarrow i_*{}^pH^0i^*F\longrightarrow0,\tag{4.1.10.1}\end{equation*}
\begin{equation*}0\longrightarrow i_*{}^pH^0i^!F\longrightarrow F\longrightarrow j_*j^*F\longrightarrow i_*{}^pH^1i^!F\longrightarrow0.\tag{4.1.10.2}\end{equation*}
断言(i)是4.1.3的特例。三角 $(j_!j^*F,F,i_*i^*F)$ 和 $(i_*i^!F,F,j_*j^*F)$ 的\term{perverse-cohomology}{反常上同调}长\term{exact-sequence}{正合列}，给出（4.1.10.1）、（4.1.10.2），并给出当 $i\ne-1,0$ 时 ${}^pH^ii_*i^*F=i_*{}^pH^ii^*F=0$，以及当 $i\ne0,1$ 时 ${}^pH^ii_*i^!F=i_*{}^pH^ii^!F=0$，因而得到(ii)。

\par\medskip\noindent\textbf{4.1.11。} 在（4.1.10.1）中取 $F=j_*G$，就得到 $i^*j_*G$ 的\term{perverse-cohomology}{反常上同调}\term{sheaf}{层}的\term{exact-sequence}{正合列}
\begin{equation*}0\longrightarrow i_*{}^pH^{-1}i^*j_*G\longrightarrow j_!G\longrightarrow j_*G\longrightarrow i_*{}^pH^0i^*j_*G\longrightarrow0.\tag{4.1.11.1}\end{equation*}
若对 $F=j_!G$ 应用（4.1.10.2），得到\term{exact-sequence}{正合列}
\[0\longrightarrow i_*{}^pH^0i^!j_!G\longrightarrow j_!G\longrightarrow j_*G\longrightarrow i_*{}^pH^1i^!j_!G\longrightarrow0.\]
通过 $i^*j_*G$ 与 $i^!j_!G[1]$ 之间的同构（1.4.6.4），这就是（4.1.11.1），只是箭头 $j_*G\to i_*{}^pH^0i^*j_*G$ 的符号相反。

\par\medskip\noindent\textbf{推论4.1.12。} 在4.1.10的假设下，$i^*j_{!*}G[-1]$ 与 $i^!j_{!*}G[1]$ 均为\term{perverse-sheaf}{反常层}，并且
\begin{equation*}i^*j_{!*}G[-1]={}^pH^{-1}i^*j_*G={}^pH^0i^!j_!G,\tag{4.1.12.1}\end{equation*}
\begin{equation*}i^!j_{!*}G[1]={}^pH^0i^*j_*G={}^pH^1i^!j_!G.\tag{4.1.12.2}\end{equation*}
\SourcePage{114}
\term{exact-sequence}{正合列}（4.1.11.1）给出两个短\term{exact-sequence}{正合列}
\begin{equation*}0\longrightarrow i_*{}^pH^{-1}i^*j_*G\longrightarrow j_!G\longrightarrow j_{!*}G\longrightarrow0,\tag{4.1.12.3}\end{equation*}
\begin{equation*}0\longrightarrow j_{!*}G\longrightarrow j_*G\longrightarrow i_*{}^pH^0i^*j_*G\longrightarrow0.\tag{4.1.12.4}\end{equation*}
分别对（4.1.12.3）、（4.1.12.4）应用 $i^*$、$i^!$，并利用 $i^*j_!=i^!j_*=0$，便得到（4.1.12.1）、（4.1.12.2），从而得到所需的反常性。

\par\medskip\noindent\textbf{注4.1.13。} 函子性的三角 $K\mapsto(j_!j^*K,K,i_*i^*K)$ 可以细化为函子
\[T:\D_c^b(X,\overline\Q_\ell)\longrightarrow\D_c^bF(X,\overline\Q_\ell),\]
使得按照3.1.7的记号，$GTK$ 为只在 $-1$、$0$ 次非零的\term{complex}{复形} $j_!j^*K\to K$。有 $\omega TK\xrightarrow\sim i_*i^*K$，而且可以选择这个同构，使
\[K=G^0TK\longrightarrow\omega TK\longrightarrow i_*i^*K\]
就是\term{adjunction}{伴随}箭头。若 $F$ 为\term{perverse-sheaf}{反常层}，则 $TF$ 是一个带愚\term{filtration}{滤过}的\term{complex}{复形}；于是 $i_*i^*F$ 与以下\term{perverse-sheaf}{反常层}\term{complex}{复形}在 $\D_c^b(X,\overline\Q_\ell)$ 中定义的对象同构：
\[i_*i^*F=[j_!j^*F\longrightarrow F]\qquad\text{（次数为 $-1$ 和 $0$）}.\]
由3.1.14从这个同构导出的四项\term{exact-sequence}{正合列}，与（4.1.10.1）相同，只是连接映射的符号相反。

同样，$K\mapsto(i_*i^!K,K,j_*j^*K)$ 可细化为 $K\mapsto TK$，其中 $GTK$ 为只在 $0$、$1$ 次非零的\term{complex}{复形} $K\to j_*j^*K$，且 $\omega TK=i_*i^!K$。因而，当 $F$ 为\term{perverse-sheaf}{反常层}时有同构
\[i_*i^!F=[F\longrightarrow j_*j^*F]\qquad\text{（次数为 $0$ 和 $1$）},\]
由此得到（4.1.10.2），至多相差符号。

\section*{4.2. 正合性与伴随}
\addcontentsline{toc}{section}{4.2. 正合性与伴随}
\par\medskip\noindent\textbf{4.2.1。} 设 $T:\mathscr D_1\to\mathscr D_2$ 为两个带 \term{t-structure}{$t$-结构}的\term{triangulated-category}{三角范畴}之间的\term{exact-functor}{正合函子}。若
\[T\mathscr D_1^{\le0}\subset\mathscr D_2^{\le b}\quad(b\ne\infty),\qquad T\mathscr D_1^{\ge0}\subset\mathscr D_2^{\ge a}\quad(a\ne-\infty),\]
就称 $T$ 的\term{cohomological-amplitude}{上同调幅度}为 $[a,b]$（$-\infty\le a\le b\le\infty$）；也就是说，$T[b]$ 右 \term{t-exact}{$t$-正合}（当 $b\ne\infty$），$T[a]$ 左 \term{t-exact}{$t$-正合}（当 $a\ne-\infty$）。若 $T$ 的\term{cohomological-amplitude}{上同调幅度}有界，这\term{equivalence}{等价}于要求：对 $A\in\mathscr C_1=\mathscr D_1^{\le0}\cap\mathscr D_1^{\ge0}$，当 $i\notin[a,b]$ 时 $H^iTA=0$。

\SourcePage{115}
\par\medskip\noindent\textbf{4.2.2。} 设 $f:X\to Y$ 是域 $k$ 上分离\term{finite-type}{有限型}\term{scheme}{概形}之间的态射。本节估计与 $f$ 相联系、连接 $\D_c^b(X,\overline\Q_\ell)$ 和 $\D_c^b(Y,\overline\Q_\ell)$ 的通常函子，相对于\term{perversity}{反常度} $p_{1/2}$ 的 \term{t-structure}{$t$-结构}的\term{cohomological-amplitude}{上同调幅度}。

\par\medskip\noindent\textbf{4.2.3。} 若 $Y$ 有\term{affine}{仿射}开覆盖 $\{V\}$，使每个 $f^{-1}V$ 都能被 $d+1$ 个\term{affine}{仿射}开集覆盖，则 $f_*$ 的\term{cohomological-amplitude}{上同调幅度} $\le d$，而 $f_!$ 的\term{cohomological-amplitude}{上同调幅度} $\ge-d$。

这两个命题互为\term{duality}{对偶}，且在 $Y$ 上是局部的，故可假设 $Y$ \term{affine}{仿射}，$X$ 可被 $d+1$ 个\term{affine}{仿射}开集覆盖。于是，$f_*$ 的断言由4.1.1及开覆盖的Leray\term{spectral-sequence}{谱序列}得到。

\par\medskip\noindent\textbf{4.2.4。} 若 $f$ 的纤维维数 $\le d$，则\term{cohomological-amplitude}{上同调幅度}有如下估计：
\[\begin{array}{c@{\qquad}c}f_!:\ \le d&f^!:\ \ge-d\\f^*:\ \le d&f_*:\ \ge-d.\end{array}\]
\term{duality}{对偶}交换每一条对角线上的两个断言；对\term{adjoint-functor}{伴随函子}对 $(f_![d],\allowbreak f^![-d])$、$(f^*[d],\allowbreak f_*[-d])$ 应用1.3.17(iii)，可知每一行的两个断言\term{equivalence}{等价}。因此只需考虑 $f^*$；它的估计直接由（4.0.1）及对任意 $x\in X$ 成立的不等式 $\dim(x)\le\dim f(x)+d$ 得到。

\noindent\textbf{注。} 同一结果在 $\Z_\ell$-\term{cohomology}{上同调}中对 \term{t-structure}{$t$-结构} $p_{1/2}$ 和 $p_{1/2}^{+}$ 也成立。只要同时处理这两个 \term{t-structure}{$t$-结构}，上述证明仍然适用。4.2.3也一样。

应用1.3.17(iii)，就得到连接 $X$、$Y$ 上\term{perverse-sheaf}{反常层}的以下\term{adjoint-functor}{伴随函子}对：
\[({}^pH^df_!,{}^pH^{-d}f^!)\qquad\text{和}\qquad({}^pH^df^*,{}^pH^{-d}f_*).\]
两个有用的特例是：

$f$ \term{proper}{固有}：$f_!=f_*$ 的\term{cohomological-amplitude}{上同调幅度}为 $[-d,d]$。

$f$ \term{smooth}{光滑}且相对维数为 $d$：有 $f^!=f^*[2d](d)$，因而
\SourcePage{116}
$f^*[d]=f^![-d](-d)$ 是 \term{t-exact}{$t$-正合}的，并在\term{perverse-sheaf}{反常层}范畴上给出三个依次\term{adjunction}{伴随}的函子
\[\bigl({}^pH^df_!(d),\ {}^pH^df^*,\ {}^pH^{-d}f_*\bigr).\]

\par\medskip\noindent\textbf{命题4.2.5。} 若 $f$ \term{smooth}{光滑}、相对维数为 $d$，且几何纤维连通（按定义，因而也非空），则 \term{t-exact}{$t$-正合}函子 $f^*[d]$ 诱导从 $Y$ 上\term{perverse-sheaf}{反常层}到 $X$ 上\term{perverse-sheaf}{反常层}的\term{fully-faithful}{全忠实}函子。

证明中会用到通常的\term{sheaf}{层}函子 $f_*$，记作 ${}^0f_*$。

对 $K,L\in\D_c^b(Y,\overline\Q_\ell)$，由于 $f$ \term{smooth}{光滑}，有
\begin{equation*}f^*R\mathcal Hom(K,L)\xrightarrow\sim R\mathcal Hom(f^*K,f^*L).\tag{4.2.5.1}\end{equation*}
对一般的 $f$，相应的公式则为
\[f^!R\mathcal Hom(K,L)\xrightarrow\sim R\mathcal Hom(f^*K,f^!L).\]
当 $K,L$ 反常时，$R\mathcal Hom(K,L)\in\D_c^{\ge0}$；\[R\mathcal Hom(f^*K,f^*L)=R\mathcal Hom(f^*K[d],f^*L[d])\] 也如此。因此，对（4.2.5.1）应用 ${}^0f_*H^0$，得到
\begin{equation*}{}^0f_*f^*H^0R\mathcal Hom(K,L)\longrightarrow{}^0f_*H^0R\mathcal Hom(f^*K[d],f^*L[d]).\tag{4.2.5.2}\end{equation*}
暂且承认，对 $Y$ 上的\term{sheaf}{层} $H$，有
\begin{equation*}H\xrightarrow\sim{}^0f_*f^*H.\tag{4.2.5.3}\end{equation*}
对（4.2.5.2）取整体截面，就得到
\[\Gamma H^0R\mathcal Hom(K,L)\xrightarrow\sim\Gamma H^0R\mathcal Hom(f^*K[d],f^*L[d]).\]
两边的 $R\mathcal Hom$ 都在 $\D_c^{\ge0}$ 中，所以
\[\Gamma H^0R\mathcal Hom=H^0R\Gamma R\mathcal Hom=H^0\RHom=\Hom,\]
命题随之得证。

（4.2.5.3）的证明可归结为集合\term{sheaf}{层}的情形。态射 $f$ 局部有截面，因此可假设 $f$ 有截面 $e$。它给出 $H\to f_*f^*H$ 的一个收缩；只须再验证，$f_*f^*H$ 的两个局部截面，若作为 $f^*H$ 的截面在 $e$ 上相同，就必相等。此事可逐纤维验证，因而归结为 $Y$ 为一点的情形，留给读者。

\SourcePage{117}
\par\medskip\noindent\textbf{4.2.6。} 设 $u^*$ 是从\term{abelian-category}{阿贝尔范畴} $\mathscr A$ 到\term{abelian-category}{阿贝尔范畴} $\mathscr B$ 的\term{fully-faithful}{全忠实}\term{exact-functor}{正合函子}，具有左、右\term{adjunction}{伴随} $u_!$、$u_*$。下列条件\term{equivalence}{等价}：
\begin{enumerate}[label=(\alph*)]
\item $u^*$ 将 $\mathscr A$ 等同于 $\mathscr B$ 的一个\term{thick-subcategory}{厚子范畴}（即对\term{subquotient}{子商}封闭的子范畴）；
\item 对 $B\in\mathscr B$，\term{adjunction-morphism}{伴随态射} $u_!u^*B\to B$ 为单态射；
\end{enumerate}
\noindent(b$'$) 对 $B\in\mathscr B$，\term{adjunction-morphism}{伴随态射} $B\to u_*u^*B$ 为满态射。\footnote{译注：本段原文的函子方向、复合次序和左右\term{adjunction}{伴随}记号之间存在不一致；这里保留原文记号。按开头给定的方向，(b)、(b$'$)中的\term{adjunction-morphism}{伴随态射}应分别为 $u^*u_*B\to B$、$B\to u^*u_!B$，而下面示例中的左右\term{adjunction}{伴随}记号相应互换。}

这些条件并非总能满足。例如，令 $\mathscr B$ 为\term{abelian-category}{阿贝尔范畴} $\mathscr C$ 中的箭头 $X\to Y$ 构成的范畴，$\mathscr A$ 为其中由同构 $X\xrightarrow\sim Y$ 构成的子范畴（它与 $\mathscr C$ \term{equivalence}{等价}），$u^*$ 为包含函子，则
\[u_!(X\to Y)=(X\xrightarrow\sim X),\qquad u_*(X\to Y)=(Y\xrightarrow\sim Y).\]
若这些条件成立，并利用 $u^*$ 将 $\mathscr A$ 看作 $\mathscr B$ 的全子范畴，则 $\mathscr B$ 的每个对象 $B$ 都有一个属于 $\mathscr A$ 的最大子对象，即 $u_!u^*B$，以及一个属于 $\mathscr A$ 的最大商，即 $u_*u^*B$。下列引理留给读者验证。

\par\medskip\noindent\textbf{引理4.2.6.1。} 若 $\mathscr A$ 的每个对象都具有有限长度，并且对 $\mathscr A$ 的每个\term{simple-object}{单对象} $S$，$u^*S$ 仍为\term{simple-object}{单对象}，则4.2.6的条件(a)、(b)、(b$'$)成立。

判别准则4.2.6.1适用于4.2.5的情形：我们有依次\term{adjunction}{伴随}的函子
\[\bigl({}^pH^df_!(d),{}^pH^df^*,{}^pH^{-d}f_*\bigr),\]
“每个对象都具有有限长度”正是4.3.1(i)。\term{simple-object}{单对象} $S=j_{!*}L[\dim V]$ 在 $f^*[d]$ 下的像为
\[
\begin{tikzcd}[ampersand replacement=\&]
f^{-1}(V)\arrow[r,"{\tilde j}"]\arrow[d,"{\tilde f}"']\&X\arrow[d,"{f}"]\\
V\arrow[r,"{j}"']\&Y
\end{tikzcd}
\]
\[\tilde j_{!*}\bigl(\tilde f^*L[\dim f^{-1}V]\bigr),\]
仍然是\term{simple-object}{单对象}。事实上，由\term{smooth}{光滑}\term{base-change}{基变换}定理，以及由2.2.4导出的、形如2.1.11的 $j_{!*}$ 描述，可知 $j_{!*}$ 与沿 $f$ 的\term{base-change}{基变换}交换。另一方面，因为 $f$ \term{smooth}{光滑}且纤维连通，$V$ 上的不可约\term{local-system}{局部系统}拉回到 $f^{-1}V$ 后仍然不可约。

\SourcePage{118}
\par\medskip\noindent\textbf{推论4.2.6.2。} 在4.2.5的假设下，从 $Y$ 上\term{perverse-sheaf}{反常层}到 $X$ 上\term{perverse-sheaf}{反常层}的函子 $f^*[d]$ 满足4.2.6中\term{equivalence}{等价}的条件(a)、(b)、(b$'$)。

对 $X$ 上的\term{perverse-sheaf}{反常层} $F$，这个推论使我们可以讨论 $F$ 中通过\term{inverse-image}{逆像}来自 $Y$ 的最大子\term{sheaf}{层}，以及具有同样性质的最大商\term{sheaf}{层}。这个子\term{sheaf}{层}是 ${}^pH^{-d}f_*F$ 的 $f^*[d]$；这个商\term{sheaf}{层}是 ${}^pH^df_!F(d)$ 的 $f^*[d]$。

上述4.2.6.2的证明简短，但也有缺点：它不能原样用于 $p_{1/2}$ 以外的\term{perversity}{反常度}，也不能原样用于 $\Z_\ell$-\term{sheaf}{层}。

\par\medskip\noindent\textbf{4.2.7。} 对域 $k$ 上的\term{finite-type}{有限型}\term{scheme}{概形} $X_1,X_2$，以 $\boxtimes$ 表示\term{external-tensor-product}{外张量积}
\[\D(X_1)\times\D(X_2)\longrightarrow\D(X_1\times X_2),\qquad(K,L)\longmapsto\operatorname{pr}_1^*K\overset L\otimes\operatorname{pr}_2^*L.\]
可以采用不同的框架：$\Z/\ell$ 系数；$\Z/\ell^n$ 系数，但限于 $\D^-$ 或 $\D_{ctf}^b$；对适当的 $k$（参见2.2），在 $\D_c^b$ 中使用 $\Z_\ell$ 或 $\Q_\ell$ 系数；等等。在所有这些框架中，\term{external-tensor-product}{外张量积}均与通常运算交换。

\noindent(a) 对 $f_1:X_1\to Y_1$、$f_2:X_2\to Y_2$，有
\[f_{1*}K\boxtimes f_{2*}L\xrightarrow\sim(f_1\times f_2)_*(K\boxtimes L).\]
对 $(f_1\times f_2)_!$、$(f_1\times f_2)^*$、$(f_1\times f_2)^!$ 也有类似结论。特别地，$X_1\times X_2$ 的\term{dualizing-complex}{对偶化复形}是 $X_1$、$X_2$ 的\term{dualizing-complex}{对偶化复形}的\term{external-tensor-product}{外张量积}。

\noindent(b) 对 $X_1$ 上的 $K_1,L_1$ 及 $X_2$ 上的 $K_2,L_2$，若 $K_1,K_2\in\D_c^b$，则
\[R\mathcal Hom(K_1,L_1)\boxtimes R\mathcal Hom(K_2,L_2)\xrightarrow\sim R\mathcal Hom(K_1\boxtimes K_2,L_1\boxtimes L_2).\]
特别地，$K\boxtimes L$ 的Verdier\term{duality}{对偶} $D(K\boxtimes L)$ 是 $K,L$ 各自Verdier\term{duality}{对偶}的\term{external-tensor-product}{外张量积}。将 $R\mathcal Hom$ 换成 $\otimes$，显然也有类似结论。

通过标准论证（参见[3]第1节），这些命题都可归结为(a)中关于\term{direct-image}{正像}的断言。将 $(f_1,f_2)$ 分解为 $(f_1,\id)\circ(\id,f_2)$，就可以假设 $f_1$ 或 $f_2$ 是恒等映射。此外还可归结为 $\Z/\ell$ 系数的情形。在此情形下，若 $f_1$ 为恒等映射，
\SourcePage{119}
则结论可由 $X_1\to\Spec(k)$ 对任意系数\term{sheaf}{层}的普遍\term{locally-acyclic}{局部无环}性推出（[3] 2.16）。

\par\medskip\noindent\textbf{命题4.2.8。} 对 $k$ 上的\term{finite-type}{有限型}\term{scheme}{概形} $X_1,X_2$，双函子 $\boxtimes$ 是 \term{t-exact}{$t$-正合}的。

设 $K_i\in\D_c^b(X_i)$（$i=1,2$），取 $X_i$ 的\term{stratification}{分层} $T_i$，使\term{stratum}{层片}等维、其既约结构在 $\bar k$ 上\term{smooth}{光滑}，而且对任一\term{stratum}{层片}的包含映射 $i_i:S\hookrightarrow X_i$，各 $H^ni_i^*K_i$、$H^ni_i^!K_i$ 均\term{smooth}{光滑}。于是，$K_i\in\D^{\le0}$（分别为 $\D^{\ge0}$）的充要条件是，对每个\term{stratum}{层片} $i:S\hookrightarrow X_i$，均有
\[H^ni^*K_i=0\quad(n>-\dim S)\qquad\text{（分别为 }H^ni^!K_i=0\quad(n<-\dim S)\text{）}.\]
将这个准则应用于 $K_1\boxtimes K_2$ 及 $X_1\times X_2$ 的\term{stratification}{分层} $T_1\times T_2$，再利用4.2.7，就得到4.2.8。这里，所用\term{perversity}{反常度}函数的线性性质是关键。

\par\medskip\noindent\textbf{4.2.9。} 同样而且更简单地，设 $K/k$ 为域\term{extension}{扩张}，$X$ 为 $k$ 上\term{finite-type}{有限型}\term{scheme}{概形}，$X_K:=X\times_{\Spec k}\Spec K$，$\varepsilon:X_K\to X$ 为投影，则函子 $\varepsilon^*$ 与通常运算交换，并且 \term{t-exact}{$t$-正合}。这对任意\term{perversity}{反常度}都成立，也适用于 $\Z_\ell$、$\Z/\ell^n$ 等系数的\term{cohomology}{上同调}。

\section*{4.3. 单对象}
\addcontentsline{toc}{section}{4.3. 单对象}
\par\medskip\noindent\textbf{定理4.3.1。} (i) $X$ 上\term{perverse-sheaf}{反常层}的范畴既\term{artinian}{阿廷}又\term{noetherian}{诺特}：每个对象都具有有限长度。

\noindent(ii) 设 $j:V\hookrightarrow X$ 为一个不可约子簇的包含映射，$V$ 在 $\bar k$ 上的既约结构\term{smooth}{光滑}，$L$ 为 $V$ 上的不可约\term{smooth}{光滑} $\overline\Q_\ell$-\term{sheaf}{层}（对应于 $\pi_1(V)$ 的一个不可约 \term{l-adic}{$\ell$-进}表示），则 $j_{!*}(L[\dim V])$ 是 $X$ 上的单\term{perverse-sheaf}{反常层}。所有单\term{perverse-sheaf}{反常层}均可如此得到。

将标量从 $k$ \term{extension}{扩张}到它的完美闭包，不改变\term{etale}{平展}拓扑。因此，可设 $k$ 完美。这时，“在 $\bar k$ 上的既约结构\term{smooth}{光滑}”可以替换为“\term{smooth}{光滑}”。

\par\medskip\noindent\textbf{引理4.3.2。} 设 $X$ 不可约且\term{smooth}{光滑}，$L$ 为 $X$ 上的\term{smooth}{光滑} $\overline\Q_\ell$-\term{sheaf}{层}。对任一非空Zariski开集 $j:U\hookrightarrow X$，\term{perverse-sheaf}{反常层} $F:=L[\dim X]$ 都满足 $F=j_{!*}j^*F$。

\SourcePage{120}
须验证（参见2.1.9），对每个 $x\in X-U$，有
\[H^ii_x^*F=0\quad(i\ge-\dim x),\qquad H^ii_x^!F=0\quad(i\le-\dim x).\]
这里 $\dim x<\dim X$；所需的第一个消失结论在
\[i>-\dim X=-\dim x-(\dim X-\dim x)\]
时成立；由\term{cohomology}{上同调}\term{pure}{纯}性定理，第二个消失结论在
\[i<-\dim X+2\operatorname{codim}x=-\dim x+(\dim X-\dim x)\]
时成立。

\par\medskip\noindent\textbf{引理4.3.3。} 设 $X$ 不可约且\term{smooth}{光滑}，$L$ 为 $X$ 上的不可约\term{smooth}{光滑} $\overline\Q_\ell$-\term{sheaf}{层}，则\term{perverse-sheaf}{反常层} $F:=L[\dim X]$ 为\term{simple-object}{单对象}。

设 $G\subset F$。存在 $X$ 的非空Zariski开集 $j:U\hookrightarrow X$，使 $G$ 在 $U$ 上形如 $M[\dim X]$，其中 $M$ 为 $L$ 的\term{smooth}{光滑}子\term{sheaf}{层}。因为 $X$ 正规，$L$ 的\term{monodromy}{单值作用}限制到 $U$ 后不变，故 $L|U$ 仍不可约。因此 $M=0$ 或 $M=L$。若 $M=0$，则 $G$ 的\term{support}{支撑}包含在闭集 $F=X-U$ 中；若 $M=L$，则 $F/G$ 的\term{support}{支撑}包含在该闭集中。由4.3.2，$F=j_{!*}j^*F$，故 $F$ 没有\term{support}{支撑}在该闭集中的非零子对象或商对象（1.4.2.5）。于是分别有 $G=0$ 或 $G=F$，4.3.3得证。

\par\medskip\noindent\textbf{4.3.4。} 下面证明4.3.1。对开浸入 $V\hookrightarrow\bar V$ 和闭浸入 $\bar V\hookrightarrow X$ 分别应用（1.4.2.5），再利用4.3.3，就知(ii)中的\term{perverse-sheaf}{反常层}都是\term{simple-object}{单对象}。为完成证明，还须验证 $X$ 上的每个\term{perverse-sheaf}{反常层}都是这些单\term{perverse-sheaf}{反常层}的逐次\term{extension}{扩张}。对 $X$ 作\term{noetherian}{诺特}归纳，因而可假设：对\term{support}{支撑} $Y\ne X$ 的\term{perverse-sheaf}{反常层}，这个结论成立。$X$ 上任一\term{perverse-sheaf}{反常层} $F$，在某个适当的不可约\term{smooth}{光滑}开集 $j:U\hookrightarrow X$ 上都形如 $L[\dim U]$，其中 $L$ 为\term{smooth}{光滑} $\overline\Q_\ell$-\term{sheaf}{层}。因此（参见4.1.11），$F$ 可分解为形如 $j_{!*}L'[\dim U]$ 的\term{sheaf}{层}及\term{support}{支撑}在 $U$ 的补集上的\term{perverse-sheaf}{反常层}，其中 $L'$ 遍历 $L$ 的不可约\term{subquotient}{子商}。结论由此得出。

\par\medskip\noindent\textbf{注4.3.5（参见4.1.9）。} 在 $\C$ 上，已知完整 $\mathscr D$-模的范畴同样\term{artinian}{阿廷}且\term{noetherian}{诺特}。这里可以是在\term{smooth}{光滑}代数簇上的代数\term{holonomic-module}{完整模}，也可以是在紧解析簇上、按解析意义定义的\term{holonomic-module}{完整模}。

\SourcePage{121}
\section*{4.4. 消失轮换（上界估计）}
\addcontentsline{toc}{section}{4.4. 消失轮换（上界估计）}
本节给出关于\term{nearby-cycles}{邻近轮换}函子正合性的一个部分结果，它是SGA 7 I 4.2的加强形式。\term{nearby-cycles}{邻近轮换}函子 $R\psi_{\bar\eta}$ 与\term{duality}{对偶}交换，这一事实确实成立，但本讲义不予证明。若结合这一事实，就可推出 $R\psi_{\bar\eta}$ 的 \term{t-exact}{$t$-正合}性。这里给出的部分结果已足以证明4.5.5，而4.5.5是我们证明\term{pure}{纯}性定理5.3.4的一个关键环节。

\par\medskip\noindent\textbf{4.4.1。} 设 $(S,\eta,s)$ 为Hensel\term{trait}{特质}，$k(\bar\eta)$ 为 $k(\eta)$ 的可分闭包，$(\bar S,\bar\eta,\bar s)$ 为 $S$ 在 $k(\bar\eta)$ 中的正规化。这里 $\bar S$ 是非离散赋值环的谱，值群为 $\Q$。设 $a:X\to S$ 为 $S$ 上的\term{finite-type}{有限型}\term{scheme}{概形}，$\bar X=X\times_S\bar S$，$X_{\bar\eta}$、$X_{\bar s}$ 为 $\bar X/\bar S$ 的两个纤维：
\[
\begin{tikzcd}[ampersand replacement=\&,column sep=large]
X_{\bar s}\arrow[r,hook,"{i}"]\arrow[d,"{a}"']\&\bar X\arrow[d,"{a}"]\&X_{\bar\eta}\arrow[l,hook',"{j}"']\arrow[d,"{a}"]\\
\bar s\arrow[r,hook,"{i}"']\&\bar S\&\bar\eta\arrow[l,hook',"{j}"]
\end{tikzcd}
\]
\term{vanishing-cycles}{消失轮换}理论中的函子 $\psi_{\bar\eta}$，就是
\[i^*j_*:\D(X_{\bar\eta})\longrightarrow\D(X_{\bar s}).\]
我们的目标是证明：对于 $X_{\bar\eta}/\bar\eta$ 和 $X_{\bar s}/\bar s$ 的\term{perversity}{反常度} $p_{1/2}$，它将 ${}^p\D_c^{\le0}$ 映入 ${}^p\D_c^{\le0}$。

我们最关心的是 $\overline\Q_\ell$-\term{cohomology}{上同调}。利用以下事实：$R\psi$ 与系数环的\term{extension}{扩张}和限制运算交换；这些运算将相应的 ${}^p\D_c^{\le0}$ 彼此映入；并且对 $K\in\D_c^b(X,\Z_\ell)$，若 $K\overset L\otimes\Z/\ell\in{}^p\D_c^{\le0}(X,\Z/\ell)$，则 $K\in{}^p\D_c^{\le0}(X,\Z_\ell)$。于是可依次将 $\overline\Q_\ell$-\term{cohomology}{上同调}归结为 $E_\lambda$-\term{cohomology}{上同调}（$E_\lambda$ 为 $\Q_\ell$ 的有限\term{extension}{扩张}）、$\Q_\ell$-\term{cohomology}{上同调}、$\Z_\ell$-\term{cohomology}{上同调}，最后归结为 $\Z/\ell$-\term{cohomology}{上同调}。在 $\Z/\ell$-\term{cohomology}{上同调}中工作的好处，是无需使用 $\Z_\ell$-\term{sheaf}{层}的运算形式。函子 $\psi_{\bar\eta}$ 保持\term{constructible}{可构造}性（SGA $4\frac12$，Th. finitude，3.2）。不过，若对域上\term{finite-type}{有限型}\term{scheme}{概形} $Y$，仿照M. Artin（SGA 4 XIV 2）直接定义 ${}^p\D^{\le0}(Y,\Z/\ell)$，就不必预先知道这个事实：它是由满足以下条件的\term{complex}{复形} $K$ 构成的范畴，即对每个点 $y$（不必为闭点）及其包含映射 $i_y:y\hookrightarrow Y$，
\SourcePage{122}
都有 $H^ii_y^*K=0$（$i>-\dim\overline{\{y\}}$）。有
\[\D_c(Y,\Z/\ell)\cap{}^p\D^{\le0}(Y,\Z/\ell)={}^p\D_c^{\le0}(Y,\Z/\ell).\]
若将\term{sheaf}{层} $F$ 写成其\term{constructible}{可构造}子\term{sheaf}{层} $F_\alpha$ 的归纳极限，则 $F[i]\in{}^p\D^{\le0}(Y,\Z/\ell)$，当且仅当每个 $F_\alpha[i]\in{}^p\D_c^{\le0}(Y,\Z/\ell)$。

\par\medskip\noindent\textbf{命题4.4.2。} $\psi_{\bar\eta}$ 将 ${}^p\D^{\le0}(X_{\bar\eta},\Z/\ell^n)$ 映入 ${}^p\D^{\le0}(X_{\bar s},\Z/\ell^n)$。

通过分解并取归纳极限，只须证明：若 $X_{\bar\eta}$ 上的\term{constructible}{可构造}\term{sheaf}{层} $F$ 的\term{support}{支撑}维数 $\le N$，则对 $X_{\bar s}$ 上位于 $x$ 的任一几何点 $\bar x$，有
\[(H^i\psi_{\bar\eta}F)_{\bar x}=0\qquad(i>N-\dim\overline{\{x\}}).\]

\par\medskip\noindent\textbf{引理4.4.3。} 当 $i>N$ 时，$H^i\psi_{\bar\eta}F=0$。

设 $x$ 为 $X_{\bar s}$ 的几何点。事实上，只考虑 $x\in X(k(\bar s))$ 就够了。将 $\bar X$ 在 $x$ 的严格Hensel化 $\bar X_{(x)}$ 写成 $\bar X$ 中 $x$ 的\term{affine}{仿射}\term{etale}{平展}邻域 $U_\alpha$ 的射影极限。有
\[(H^i\psi_{\bar\eta}F)_x=\varinjlim_\alpha H^i(U_{\alpha\bar\eta},F).\]
在右侧 $H^i$ 的括号内，可将 $U_{\alpha\bar\eta}$ 换成 $F$ 的\term{support}{支撑}在 $U_{\alpha\bar\eta}$ 中的\term{inverse-image}{逆像}；再应用\term{affine}{仿射}\term{scheme}{概形}的\term{cohomology}{上同调}维数定理SGA 4 XIV 3.2，便得结论。

\par\medskip\noindent\textbf{4.4.4。} 证明一个比4.4.2略一般的命题会更方便。设 $S^h$ 为 $S$ 的严格Hensel化，其分式域为 $k(\eta^h)\subset k(\bar\eta)$；设 $k(\eta')$ 为 $k(\bar\eta)$ 中 $k(\eta^h)$ 的任一\term{extension}{扩张}，包含某个一致化元的所有次数为 $\ell$ 的幂的根。在4.4.1中用 $k(\eta')$ 替换 $k(\bar\eta)$，就定义了 $S'$、$X'$，以及函子
\[\psi_{\eta'}:\D(X_{\eta'})\longrightarrow\D(X_{\bar s}).\]
我们将对此函子证明4.4.2的类似结论。

Galois群 $Q=\Gal(k(\bar\eta)/k(\eta'))$ 是阶数与 $\ell$ 互素的有限群的射影极限。因此，“取 $Q$-不变量”是\term{exact-functor}{正合函子}。对 $K'\in\D(X_{\eta'})$，设其在 $\D(X_{\bar\eta})$ 中的\term{inverse-image}{逆像}为 $K$，则
\begin{equation*}H^i\psi_{\eta'}K'=(H^i\psi_{\bar\eta}K)^Q.\tag{4.4.4.1}\end{equation*}
验证留给读者。我们并不需要\term{derived-category}{导出范畴}中的相应命题。

\SourcePage{123}
由原命题及（4.4.4.1），即可得到4.4.2对 $\psi_{\eta'}$ 的类似结论。同样，由4.4.3及（4.4.4.1），得到4.4.3对 $\psi_{\eta'}$ 的类似结论。

设 $A$ 在 $S$ 上\term{smooth}{光滑}，$\bar s_1$ 为 $A_{\bar s}$ 的一个几何一般点，其在 $A_s$ 中的像为 $s_1$。设 $S_1$ 为 $A$ 在 $s_1$ 的Hensel化，它是一个\term{trait}{特质}。再令 $S_1^h$ 为它相对于 $\bar s_1$ 的严格Hensel化。它映到 $S^h$，并且 $S^h$ 的一致化元的\term{inverse-image}{逆像}仍为 $S_1^h$ 的一致化元。纤维积 $S_1'=S'\times_{S^h}S_1^h$ 对 $S_1$ 而言，属于上面对 $S$ 所考虑的类型。

\par\medskip\noindent\textbf{4.4.5。} 我们来证明4.4.2的4.4.4加强形式。问题在 $X$ 上是局部的，因此可以并且将假设 $X$ \term{affine}{仿射}。设 $A_S^d$ 为 $S$ 上的相对 $d$ 维\term{affine}{仿射}空间。我们将证明以下命题，4.4.2可由它形式地推出：对每个 $d$ 及每个态射 $f:X\to A_S^d$，若 $X_{\bar s}$ 的几何点 $x$ 在 $A_{\bar s}^d$ 中的像为一般点，则
\[(H^j\psi_{\eta'}F)_x=0\qquad(j>N-d),\]
其中 $F$ 如4.4.2之后所述。

对 $A_S^d/S$ 应用4.4.4：
\[
\begin{tikzcd}[ampersand replacement=\&,column sep=large]
A_S^d\arrow[dr]\&S_1\arrow[l]\arrow[d]\&S_1^h\arrow[l]\arrow[d]\&S_1'\arrow[l]\arrow[d]\\
\&S\&S^h\arrow[l]\&S'\arrow[l]
\end{tikzcd}
\]
令 $X_1$ 为 $X$ 的“\term{localization}{局部化}” $X\times_{A_S^d}S_1$，再将基\term{scheme}{概形}从 $S_1$ 换成 $S_1'$，得到 $X_1'$：
\[
\begin{tikzcd}[ampersand replacement=\&,column sep=large]
X'\arrow[d]\arrow[dr,"{f}"']\&\&X_1'\arrow[ll,"{\varepsilon}"']\arrow[d]\\
S'\&A_{S'}^d\arrow[l]\&S_1'\arrow[l]
\end{tikzcd}
\]
由于态射 $\varepsilon$ 是\term{projective-limit}{逆极限}\term{etale}{平展}的，关于 $X'/S'$ 的函子 $\psi_{\eta'}$ \term{localization}{局部化}后，就是关于 $X_1'/S_1'$ 的相应函子：
\[\psi_{\eta_1'}\varepsilon^*=\varepsilon^*\psi_{\eta'}.\]
若 $F$ 为 $X_{\eta'}'$ 中维数 $\le N$ 的闭集，则其在 $(X_1')_{\eta_1'}$ 中的\term{inverse-image}{逆像}维数 $\le N-d$。同样，若 $x\in X_{\bar s}$ 且 $f(x)$ 为 $A_{\bar s}^d$ 的一般点，即 $x$ 为 $(X_1')_{\bar s_1}$ 中某一点 $x_1$ 的像，
\SourcePage{124}
则
\[\dim\overline{\{x_1\}}=\dim\overline{\{x\}}-d.\]
这个 $d$ 次的维数差，以及上面 $X'/S'$ 与 $X_1'/S_1'$ 的\term{nearby-cycles}{邻近轮换}之间的关系，就将问题归结为 $X_1'/S_1'$ 的4.4.3（以及4.4.4给出的加强形式）。

\par\medskip\noindent\textbf{注4.4.5。} 我们从\term{affine}{仿射}\term{scheme}{概形}的\term{cohomology}{上同调}维数定理推出了4.4.2。这样的表述略显人为：在SGA 4 XIV中，这个定理与4.4.2的某种形式（隐含在SGA 4 XIV 4.5的证明中），是同时对 $N$ 归纳证明的。

\section*{4.5. Betti数的估计}
\addcontentsline{toc}{section}{4.5. Betti数的估计}
本节给出\term{betti-number}{Betti数}的估计，目标是4.5.5，它将用于\term{pure}{纯}性定理5.3.4的证明。

\par\medskip\noindent\textbf{命题4.5.1。} 设 $(X_a)_{1\le a\le n}$ 为代数闭域 $k$ 上的连通\term{smooth}{光滑}射影曲线，$X$ 为它们的乘积，$K\in{}^p\D_c^{\le0}(X,\overline\Q_\ell)$。存在常数 $C$，使得对任意 $X_a$ 的连通\term{etale}{平展}覆盖 $\widetilde X_a$，设其次数为 $d_a$，以 $\widetilde X$ 表示各 $\widetilde X_a$ 的乘积，则有：
\begin{enumerate}[label=(\roman*)]
\item 对每个 $i$，$\displaystyle\dim H^i(\widetilde X,K)\le C\prod_a d_a$；
\item 对 $i>0$，$\displaystyle\dim H^i(\widetilde X,K)\le C\sup\left\{\prod_{a\in A}d_a\;\middle|\;|A|=n-i\right\}$（当 $i>n$ 时，上确界取为零）。
\end{enumerate}
只须结论对各个 $H^jK[-j]$ 成立，就对 $K$ 成立。因此，可设 $K=F[N]$，其中 $F$ 为\term{support}{支撑}维数 $\le N$ 的 $\overline\Q_\ell$-\term{sheaf}{层}。由 $\overline\Q_\ell$-\term{sheaf}{层}可归结为 $\Q_\ell$-\term{sheaf}{层}，再归结为 $\Z_\ell$-\term{sheaf}{层}（将维数换成秩）。泛系数公式，即\term{exact-sequence}{正合列}
\[
\begin{aligned}
0\longrightarrow H^i(X,K)\otimes\Z/\ell&\longrightarrow H^i(X,K\overset L\otimes\Z/\ell)\\
&\longrightarrow\operatorname{Tor}_1(H^{i+1}(X,K),\Z/\ell)\longrightarrow0,
\end{aligned}
\]
表明 $\rank H^i(X,K)\le\dim_{\Z/\ell}H^i(X,K\overset L\otimes\Z/\ell)$，从而最终归结为 $\Z/\ell$-模\term{sheaf}{层}的情形。这样就不再需要 \term{l-adic}{$\ell$-进}运算形式。

因此，设 $F$ 为 $\Z/\ell$-模\term{sheaf}{层}，$\dim\Supp F\le N$，$K=F[N]$，并对 $n$ 归纳证明4.5.1。

\SourcePage{125}
$n=0$ 的情形留给读者。

设 $n\ge1$，以 $\operatorname{pr}$ 表示 $X\to X_1$ 的投影。存在非空开集 $j_1:U_1\hookrightarrow X_1$，使 $K$ 在其上方相对于 $\operatorname{pr}$ \term{locally-acyclic}{局部无环}（SGA $4\frac12$，finitude，2.13）。取 $U_1\ne X_1$。再以 $j$ 表示 $U=\operatorname{pr}^{-1}(U_1)$ 到 $X$ 的包含映射，以 $i$ 表示其补闭集的包含映射。利用三角 $(j_!j^*K,K,i_*i^*K)$ 分解 $K$，可归结为以下两种情形：

\noindent(a) $K$ 形如 $(j_!F)[N]$，其中 $F$ 为 $U$ 上相对于 $\operatorname{pr}$ \term{locally-acyclic}{局部无环}的\term{sheaf}{层}；

\noindent(b) $K$ 的\term{support}{支撑}包含在 $X-U$ 中。

情形(b)可归结为：对某个适当闭点 $t\in X_1$，$K$ 的\term{support}{支撑}包含在 $\operatorname{pr}^{-1}(t)$ 中。令
\[Y=\operatorname{pr}^{-1}(t)\simeq\prod_{a\ne1}X_a,\qquad\widetilde Y=\prod_{a\ne1}\widetilde X_a.\]
有
\[\dim H^i(\widetilde X,K)=d_1\dim H^i(\widetilde Y,K).\]
对 $Y$ 应用归纳假设，即得4.5.1。

现设 $K$ 形如(a)。以 $\tilde j:\widetilde U\hookrightarrow\widetilde X$、$\tilde j_1:\widetilde U_1\hookrightarrow\widetilde X_1$ 表示 $U$、$U_1$ 在 $\widetilde X$、$\widetilde X_1$ 中的\term{inverse-image}{逆像}，以 $\widetilde{\operatorname{pr}}$ 表示 $\widetilde X\to\widetilde X_1$ 或 $\widetilde U\to\widetilde U_1$ 的投影。$F$ 限制到 $\widetilde U$ 后仍相对于 $\widetilde{\operatorname{pr}}$ \term{locally-acyclic}{局部无环}，故各个 $R^q\widetilde{\operatorname{pr}}_*F$ 在 $\widetilde U_1$ 上局部常值。

考虑Leray\term{spectral-sequence}{谱序列}
\[E_2^{pq}=H^p(\widetilde X_1,R^q\widetilde{\operatorname{pr}}_*K)\Longrightarrow H^{p+q}(\widetilde X,K).\]
由于 $R^q\widetilde{\operatorname{pr}}_*K=\tilde j_{1!}R^{q+N}\widetilde{\operatorname{pr}}_*F$ 是 $\widetilde U_1$ 上\term{local-system}{局部常值层}的\term{extension-by-zero}{零延拓}，又因我们取了 $U_1\ne X_1$，从而 $\widetilde U_1\ne\widetilde X_1$，故 $p\ne1,2$ 时 $E_2^{pq}=0$。只须对 $E_2^{pq}$ 验证4.5.1的估计（其中 $i=p+q$）。\term{equivalence}{等价}地，也更方便地，只须对 $E_2^{2q}$（其中 $i=q+2$）及
\[\bigl|\chi(\widetilde X_1,R^q\widetilde{\operatorname{pr}}_*K)\bigr|=\bigl|\dim E_2^{1q}-\dim E_2^{2q}\bigr|\qquad(i=q+1)\]
作此验证。

取 $U_1$ 中的闭点 $t$ 及其上方的 $\tilde t\in\widetilde U_1$，$Y$、$\widetilde Y$ 如前，并以 $K_1$ 表示 $K$ 在 $\operatorname{pr}^{-1}(t)\simeq Y$ 上的限制。若 $t$ 充分一般，则
\[\dim(\Supp F\cap\operatorname{pr}^{-1}t)\le\dim\Supp F-1,\]
故 $K_1[-1]\in{}^p\D_c^{\le0}(Y,\Z/\ell)$。
\SourcePage{126}
而 $E_2^{2q}=H_c^2(\widetilde U_1,R^q\widetilde{\operatorname{pr}}_*K)$ 由\term{monodromy}{单值作用}在
\[(R^q\widetilde{\operatorname{pr}}_*K)_{\tilde t}=H^q(\widetilde Y,K_1)\]
上的余不变量给出：
\[E_2^{2q}=H^q(\widetilde Y,K_1)_{\pi_1(\widetilde U_1,\tilde t)}(-1).\]
于是 $\dim E_2^{2q}\le\rank R^q\widetilde{\operatorname{pr}}_*K$。又因
\begin{equation*}\rank R^q\widetilde{\operatorname{pr}}_*K=\dim H^q(\widetilde Y,K_1)=\dim H^{q+1}(\widetilde Y,K_1[-1]),\tag{4.5.1.1}\end{equation*}
所需的 $E_2^{2q}$ 估计由归纳假设得到。

因为 $\widetilde X_1$ 是\term{proper}{固有}曲线 $X_1$ 的\term{etale}{平展}覆盖，也可由下述Euler--Poincaré公式知道，$\chi(\widetilde X_1,R^q\widetilde{\operatorname{pr}}_*K)$ 是投影 $X_1\times\prod_{a\ne1}\widetilde X_a\to X_1$ 的相应Euler示性数的 $d_1$ 倍。因此，我们只须在 $d_1=1$，即 $\widetilde X_1=X_1$ 的情形证明下列估计；下文均作此假设：
\[
\bigl|\chi(X_1,R^q\widetilde{\operatorname{pr}}_*K)\bigr|\le
\begin{cases}
O\!\left(\displaystyle\prod_{a\ne1}d_a\right),\\[4pt]
O\!\left(\displaystyle\sup_{\substack{A\subset[2,n]\\ |A|=n-(q+1)-1}}\prod_{a\in A}d_a\right),&q\ge0.
\end{cases}
\]
由Euler--Poincaré公式（参见[9]），有
\[\chi(X_1,R^q\widetilde{\operatorname{pr}}_*K)=\chi(U_1)\rank(R^q\widetilde{\operatorname{pr}}_*K)-\sum_{x\in X_1-U_1}\operatorname{Sw}_x(R^q\widetilde{\operatorname{pr}}_*K).\]
第一项可由秩的估计（4.5.1.1）及归纳假设处理。为了给各个\term{swan-conductor}{Swan导子} $\operatorname{Sw}_x(R^q\widetilde{\operatorname{pr}}_*K)$ 以界，我们将对 $X_1$ 在 $x$ 的Hensel化\term{trait}{特质} $(X_{1(x)},\eta,x)$，以及由 $X\to X_1$、$\widetilde X\to X_1$ \term{base-change}{基变换}得到的
\[X_{(x)}\longrightarrow X_{1(x)},\qquad\widetilde X_{(x)}\longrightarrow X_{1(x)}\]
应用\term{vanishing-cycles}{消失轮换}理论。设 $K_1$ 为 $K$ 在 $(X_{(x)})_\eta$ 上的\term{inverse-image}{逆像}。同样，$K_1[-1]\in{}^p\D_c^{\le0}(X_{(x)\eta},\Z/\ell)$。因此，在 $X_{(x)x}=\operatorname{pr}^{-1}(x)\xrightarrow\sim Y$ 上的\term{complex}{复形} $\psi_{\bar\eta}K_1[-1]$ 属于 ${}^p\D_c^{\le0}(Y,\Z/\ell)$（4.4.2）。$\widetilde X_{(x)x}=\widetilde{\operatorname{pr}}^{-1}(x)\xrightarrow\sim\widetilde Y$ 上的相应\term{complex}{复形}由它取\term{inverse-image}{逆像}得到。这些\term{complex}{复形}带有惯性群 $I=\Gal(\bar\eta/\eta)$ 的作用，特别是带有野惯性群 $P$ 的作用，并有\term{equivariant}{等变}同构
\[(R^q\widetilde{\operatorname{pr}}_*K)_{\bar\eta}=H^q(\widetilde Y,\psi_{\bar\eta}K_1)=H^{q+1}(\widetilde Y,\psi_{\bar\eta}K_1[-1]).\]
$P$ 在 $\psi_{\bar\eta}K_1$ 上的作用经过某个有限商 $Q$ 因子化。
\SourcePage{127}
事实上，只要知道 $P$ 在各个 $H^i\psi_{\bar\eta}K_1$ 上的作用具有这样的因子化就够了；不过，把 $\psi_{\bar\eta}K_1$ 看成 $\D(Y,\Z/\ell[Q])$ 中的对象，证明会更清楚。设 $\operatorname{Sw}$ 为 $Q$ 的Swan表示，它是\term{projective}{投射}的 $\Z/\ell[Q]$-模，而\term{swan-conductor}{Swan导子}为
\[\operatorname{Sw}_xR^q\widetilde{\operatorname{pr}}_*K:=\dim\Hom_Q\bigl(\operatorname{Sw},(R^q\widetilde{\operatorname{pr}}_*K)_{\bar\eta}\bigr),\]
并且
\[
\begin{aligned}
\Hom_Q\bigl(\operatorname{Sw},(R^q\widetilde{\operatorname{pr}}_*K)_{\bar\eta}\bigr)
=H^{q+1}\bigl(\widetilde Y,\mathcal Hom_Q(\operatorname{Sw},\psi_{\bar\eta}K_1[-1])\bigr).
\end{aligned}
\]
归纳假设适用于 $Y$ 和 $\mathcal Hom_Q(\operatorname{Sw},\psi_{\bar\eta}K_1[-1])$，应用它即完成证明。

\par\medskip\noindent\textbf{推论4.5.2。} 对 $K\in\D_c^b(X,\overline\Q_\ell)$，存在常数 $C$，使对任意 $\widetilde X_a$ 及任意 $i$，都有
\[\dim H^i(\widetilde X,K)\le C\prod_a d_a.\]
作平移即可归结为4.4.1(i)。\footnote{译注：原文如此；所指应为4.5.1(i)。}

\par\medskip\noindent\textbf{推论4.5.3。} 对 $K\in{}^p\D_c^{\ge0}(X,\overline\Q_\ell)$，存在常数 $C$，使对任意 $\widetilde X_a$ 及任意 $i\le0$，都有
\[\dim H^i(\widetilde X,K)\le C\sup\left\{\prod_{a\in A}d_a\;\middle|\;|A|=n-|i|\right\}.\]
这是4.5.1(ii)的\term{duality}{对偶}命题。

\par\medskip\noindent\textbf{推论4.5.4。} 设 $f:U\to X$ 拟有限且分离，$K\in{}^p\D_c^{\ge0}(U,\overline\Q_\ell)$。存在常数 $C$，使得对任意 $\widetilde X_a$，令 $\widetilde U=U\times_X\widetilde X$，都有：对每个 $i$，
\[\dim H^i(\widetilde U,K)\le C\prod_a d_a;\]
对 $i\le0$，
\[\dim H^i(\widetilde U,K)\le C\sup\left\{\prod_{a\in A}d_a\;\middle|\;|A|=n-|i|\right\}.\]
由于 $H^*(\widetilde U,K)=H^*(\widetilde X,f_*K)$，且由2.2.5，$f_*K\in{}^p\D_c^{\ge0}(X,\overline\Q_\ell)$，对它应用4.5.3即可。

\par\medskip\noindent\textbf{推论4.5.5。} 设 $X$ 为 $k$ 上\term{finite-type}{有限型}\term{scheme}{概形}，$F$ 为 $X$ 上的\term{perverse-sheaf}{反常层}，$x\in X$。存在 $x$ 的一个\term{affine}{仿射}\term{etale}{平展}邻域 $U$，以及 $U$ 的一族\term{etale}{平展}覆盖 $U_n$，其次数 $d_n$ 趋于无穷，并且当 $i<0$ 时
\[\dim H^i(U_n,F)\le o(d_n).\]
问题是局部的，因此可设 $X$ \term{affine}{仿射}。取一个拟有限态射
\SourcePage{128}
$f:X\to A^d$（例如浸入），使 $f(x)=0$。设 $h:C\to\mathbb P^1$ 为\term{smooth}{光滑}射影曲线到 $\mathbb P^1$ 的态射，$C$ 的亏格 $g>0$，并且 $h$ 在位于 $0$ 上方的某个 $c\in C$ 处\term{etale}{平展}。取包含 $c$ 的开集 $C'$，使它在 $\mathbb P^1$ 上\term{etale}{平展}，且 $h(C')\subset A^1$。设 $C_n$ 为 $C$ 的连通\term{etale}{平展}覆盖，其次数 $d_n'$ 趋于无穷（这里用到了 $g>0$）；以 $C_n'$ 表示 $C'$ 在 $C_n$ 中的\term{inverse-image}{逆像}。令 $U$、$U_n$ 分别为纤维积
\[X\times_{A^d}(C')^d,\qquad X\times_{A^d}(C_n')^d:\]
\[
\begin{tikzcd}[ampersand replacement=\&,column sep=large]
U_n\arrow[r,"{f}"]\arrow[d]\&(C_n')^d\arrow[r,hook]\arrow[d]\&C_n^d\arrow[d]\\
U\arrow[r,"{f}"]\arrow[d]\&(C')^d\arrow[r,hook]\arrow[d]\&C^d\\
X\arrow[r,"{f}"']\&A^d\&
\end{tikzcd}
\]
\term{etale}{平展}覆盖 $U_n/U$ 的次数为 $d_n=(d_n')^d$；另一方面，由4.5.4，对 $i<0$，有
\[\dim H^i(U_n,F)\le C(d_n')^{d-|i|}.\]

\par\medskip\noindent\textbf{注4.5.6。} 这个证明表明：若 $X$ 的维数为 $d$，就可选取 $U_n$，使
\[\dim H^i(U_n,F)\le O\bigl(d_n^{1-|i|/d}\bigr).\]
